%=============================================================================
\section{Spectral Framework for Pretraining Monitoring}
\label{sec:framework}
%=============================================================================

We define two primary spectral metrics that together provide a complete ``instrument panel'' for pretraining: the \emph{normalized stable rank} (SR/$d$), which measures compression progress, and the \emph{power-law exponent} ($\alpha$), which measures structural quality. Both are computable from weights alone, without training data or gradient statistics.

\subsection{Metric Definitions}
\label{sec:metrics}

\paragraph{Stable rank (SR).} For a weight matrix $W \in \R^{m \times n}$ with singular values $\sigma_1 \geq \sigma_2 \geq \cdots \geq \sigma_{\min(m,n)}$, the stable rank is:
\begin{equation}
    \mathrm{SR}(W) = \frac{\|W\|_F^2}{\|W\|_{\mathrm{op}}^2} = \frac{\sum_i \sigma_i^2}{\sigma_1^2}.
    \label{eq:stable_rank}
\end{equation}
This measures the \emph{effective dimensionality}: how many ``$\sigma_1$-sized'' directions the matrix uses. Crucially, SR requires only the Frobenius norm ($O(mn)$) and the top singular value (power iteration, $O(mn)$)---no full SVD.

\paragraph{Normalized stable rank (SR/$d$).} We average across all 2D weight layers and normalize by hidden dimension $d$:
\begin{equation}
    \mathrm{SR}/d = \frac{1}{L} \sum_{l=1}^L \frac{\mathrm{SR}(W_l)}{d},
    \label{eq:sr_norm}
\end{equation}
where $L$ is the number of measured layers. This normalization makes the metric comparable across architectures.

\paragraph{Power-law exponent ($\alpha$).} Following Martin \& Mahoney~\citep{martin_mahoney_jmlr2021}, we fit the empirical spectral density (ESD) of each layer's eigenvalues ($\lambda_i = \sigma_i^2$) to a power law $P(\lambda) \sim \lambda^{-\alpha}$ using MLE with KS-statistic-based $x_{\min}$ selection. The global $\alpha$ is the parameter-weighted average across layers:
\begin{itemize}[leftmargin=*,itemsep=0pt]
    \item $\alpha > 6$: \textbf{Random} (Marchenko-Pastur bulk, no learned structure)
    \item $\alpha \in [4, 6]$: \textbf{Bulk + Spikes} (partially trained, isolated features)
    \item $\alpha \in [2, 4]$: \textbf{Heavy-Tail} (well-trained, pervasive self-regularization)
    \item $\alpha < 2$: \textbf{Rank collapse} (over-trained, potential failure mode)
\end{itemize}

\paragraph{Supplementary metrics.} Top-10 concentration $C_{10} = \sum_{i=1}^{10}\sigma_i^2 / \sum_j \sigma_j^2$; layer-type decomposition ($\alpha_{\mathrm{attn}}$ vs.\ $\alpha_{\mathrm{mlp}}$); and $\alpha$ velocity $d\alpha/dt$.

\subsection{Information-Theoretic Foundation}
\label{sec:theory}

The connection between stable rank and entropy is not an analogy---it is a mathematical identity.

\paragraph{Stable rank as spectral entropy.}
Define the \emph{spectral probability distribution} over squared singular values:
\begin{equation}
    p_i = \frac{\sigma_i^2}{\sum_j \sigma_j^2}, \qquad \sum_i p_i = 1.
    \label{eq:spectral_prob}
\end{equation}
The R\'{e}nyi-2 entropy of this distribution is $H_2 = -\log\!\bigl(\sum_i p_i^2\bigr)$, and the \emph{participation ratio} (exponential of R\'{e}nyi-2 entropy) is:
\begin{equation}
    \mathrm{PR} = \exp(H_2) = \frac{1}{\sum_i p_i^2} = \frac{\bigl(\sum_i \sigma_i^2\bigr)^2}{\sum_i \sigma_i^4}.
    \label{eq:participation_ratio}
\end{equation}
The stable rank relates to PR via the inequality chain:
\begin{equation}
    1 \;\leq\; \mathrm{SR} \;\leq\; \mathrm{PR} = \exp(H_2) \;\leq\; \exp(H_1) \;\leq\; \mathrm{rank}(W),
    \label{eq:entropy_chain}
\end{equation}
where $H_1 = -\sum_i p_i \log p_i$ is the Shannon entropy of $\{p_i\}$. Thus $\log(\mathrm{SR})$ provides a \emph{lower bound} on the spectral Shannon entropy. Equality $\mathrm{SR} = \mathrm{PR}$ holds when the spectrum follows a geometric decay---precisely the regime trained transformers occupy.

\paragraph{Universal compression as entropy reduction.}
At initialization, Marchenko-Pastur theory gives $\mathrm{SR}_{\mathrm{init}} \approx mn/(\sqrt{m}+\sqrt{n})^2 \approx 0.4d$. The empirical convergence to SR/$d \in [0.043, 0.074]$ (Section~\ref{sec:universal}) implies:
\begin{equation}
    \Delta H_2 \;\approx\; \log\!\bigl(\mathrm{SR}_{\mathrm{final}}/\mathrm{SR}_{\mathrm{init}}\bigr) \;=\; \log(\mathrm{UCR}) \;\approx\; -2.0 \;\text{nats}.
    \label{eq:delta_entropy}
\end{equation}
This is not a metaphor: training universally removes ${\approx}2$ nats of R\'{e}nyi-2 spectral entropy per layer, regardless of model scale or architecture family (verified across 70M--6.9B; larger models that have not converged show smaller compression, consistent with training insufficiency rather than a different asymptote). The Universal Compression Ratio $\mathrm{UCR} = \mathrm{SR}_{\mathrm{final}}/\mathrm{SR}_{\mathrm{init}} \approx 0.13$ is the fraction of initial spectral dimensionality retained.

\subsection{Thermodynamic Grounding}
\label{sec:thermo}

\paragraph{SGD as Langevin dynamics.}
SGD with learning rate $\eta$, weight decay $\lambda$, and mini-batch noise on a weight matrix $W$ follows the overdamped Langevin equation~\citep{liu_tegmark2025, pvnkt_ijcai2026}:
\begin{equation}
    dW = -\eta\bigl(\nabla_W \mathcal{L} + \lambda W\bigr)\,dt + \sqrt{2\eta T_{\mathrm{eff}}}\,d\xi,
    \label{eq:langevin}
\end{equation}
where $T_{\mathrm{eff}} = \eta\,\sigma^2_{\mathrm{grad}} / (2B)$ is the effective temperature set by learning rate, gradient noise variance $\sigma^2_{\mathrm{grad}}$, and batch size $B$~\citep{liu_tegmark2025}. In the stationary regime, singular values are distributed according to:
\begin{equation}
    P(\sigma) \;\propto\; \sigma^{-\alpha}, \qquad \alpha = 1 + \frac{\lambda}{T_{\mathrm{eff}}} \cdot f(\mathrm{SNR}),
    \label{eq:alpha_temperature}
\end{equation}
where $f(\mathrm{SNR})$ encodes the signal-to-noise ratio of data-driven gradients vs.\ stochastic fluctuations~\citep{sadrtdinov2025}. This provides a mechanistic interpretation: $\alpha$ measures the balance between the data-driven ``field'' organizing the spectrum and the thermal fluctuations randomizing it.

\paragraph{$\alpha$ reversal as a phase transition.}
Define the order parameter $\varphi = (\alpha_{\mathrm{rand}} - \alpha)/(\alpha_{\mathrm{rand}} - \alpha_{\mathrm{opt}})$, mapping the spectral state onto $[0,1]$ where $\varphi = 0$ is unstructured and $\varphi = 1$ is maximally organized. Training drives the system toward $\varphi > 0$ (ordered phase) as the data-driven field $h \propto \mathrm{SNR}$ dominates thermal noise $T_{\mathrm{eff}}$. Once data is exhausted ($D/N < \tau$), the signal saturates but $T_{\mathrm{eff}}$ persists:
\begin{equation}
    \frac{d\varphi}{dt} \;\propto\; h(t) - T_{\mathrm{eff}} \cdot \varphi.
    \label{eq:order_param}
\end{equation}
When $h(t) \to 0$ (no new information from repeated data), the system relaxes back toward disorder: $\alpha$ increases. This is the $\alpha$~reversal---a continuous phase transition from an ordered (heavy-tail) to a disordered (bulk+spikes) spectral phase, driven by exhaustion of the data signal relative to optimizer temperature. The critical point occurs at $D/N \approx \tau$, which is model-size-dependent: $\tau \approx 250$--$300$ for models $\leq$1B, but grows with $N$ for larger models (Section~\ref{sec:scaling}).

\paragraph{Free energy interpretation.}
Following~\citet{sadrtdinov2025}, the training objective implicitly minimizes free energy $\mathcal{F} = U - T_{\mathrm{eff}}\,\mathcal{S}$, where $U$ is the loss landscape energy and $\mathcal{S}$ the spectral entropy. Stable rank tracks $\mathcal{S}$; $\alpha$ tracks the \emph{gradient} $\partial \mathcal{S}/\partial U$ (how efficiently entropy reduction translates to loss reduction). The heavy-tail phase ($\alpha \in [2,4]$) corresponds to the thermodynamically efficient regime where each nat of entropy removed yields maximal loss decrease~\citep{liu_ziyin_neurips2025, freedman2026}.

\subsection{Practical Computation}
\label{sec:computation}

\begin{table}[t]
\centering
\small
\caption{Computational cost of spectral metrics. All measurements are embarrassingly parallel across layers and checkpoints.}
\label{tab:cost}
\begin{tabular}{@{}lllp{5cm}@{}}
\toprule
\textbf{Metric} & \textbf{Cost per layer} & \textbf{Full SVD needed?} & \textbf{Note} \\
\midrule
SR (stable rank) & $O(mn)$ & No & Power iteration for $\sigma_1$ + Frobenius norm \\
$\alpha$ (power law) & $O(k^2 \max(m,n))$ & Partial (top-$k$) & Randomized SVD with $k\!=\!256$ \\
$C_{10}$ (concentration) & $O(k^2 \max(m,n))$ & Partial (top-10) & Same SVD as $\alpha$ \\
\bottomrule
\end{tabular}
\end{table}

For a 7B model with ${\sim}100$ weight layers: full measurement takes ${\sim}5$ minutes on a single GPU. Measured every 5{,}000 steps over 143K total $= 29$ measurements $\approx 2.4$ GPU-hours, versus ${\sim}5{,}000$ GPU-hours for training: \textbf{$<0.05\%$ overhead}.

\subsection{Key Predictions}
\label{sec:predictions}

The framework yields four testable predictions:

\noindent\textbf{P1 (Universal compression).} $\mathrm{SR}/d$ converges to $\approx 0.04$--$0.07$ regardless of scale $N$ or architecture, corresponding to $\Delta H_2 \approx -2$ nats universally.
\emph{Falsified if}: SR/$d$ varies by $>50\%$ across scales or architectures.

\noindent\textbf{P2 (Performance prediction).} SR/$d$ has stronger rank correlation with downstream performance than any single scalar from the loss curve.
\emph{Falsified if}: Spearman $|r| < 0.7$.

\noindent\textbf{P3 ($\alpha$ reversal as phase transition).} For models with $D/N < \tau$, $\alpha$ reaches a minimum then \emph{increases}, signaling the order-disorder transition when data signal is exhausted relative to $T_{\mathrm{eff}}$.
\emph{Falsified if}: $\alpha$ monotonically decreases for all models regardless of $D/N$.

\noindent\textbf{P4 (Structural Chinchilla).} For models $\leq$1B, $\alpha_{\mathrm{final}}$ follows $\alpha = \alpha_\infty + A\exp(-D/\tau N)$ with $\tau \approx 269$. For larger models, structural maturity requires substantially more tokens per parameter, suggesting $\tau$ grows with $N$.
\emph{Falsified if}: no relationship between $D/N$ and $\alpha_{\mathrm{final}}$, or large models achieve low $\alpha$ with small $D/N$.
